
\section{Result status and remaining questions}\label{sec:status}
\textbf{PROVED} denotes a complete supplied argument, with finite checks where useful. It does not denote independent peer review, formal proof verification or certified priority.

\begin{table}[ht]
\centering\small
\begin{tabularx}{\textwidth}{@{}lYl@{}}
\toprule Item & Scope & Status\\\midrule
Theorems 1--6 & Original geometry, honest interval and comparison bounds & PROVED\\
Proposition~\ref{prop:diagonal_minimax} & Classical diagonal-information optimum & PROVED\\
Theorems~\ref{thm:scan}--\ref{thm:harmonic} & Initial full-data scan and harmonic obstruction & PROVED\\
Theorem~\ref{thm:sparse_harmonic} & Sharp cubic sparse harmonic lower bound & PROVED\\
Theorem~\ref{thm:full_minimax} & Matching full-data zero-point and worst-case lengths & PROVED\\
Proposition~\ref{prop:mse_sharp} & Full-data squared-error minimax rates & PROVED\\
Proposition~\ref{prop:truncated} & Certified finite-scan search tradeoff & PROVED\\
Theorem~\ref{thm:low_degree} & Restricted polynomial-estimation obstruction & PROVED\\
Proposition~\ref{prop:ci_detection} & Interval-to-testing implication & PROVED\\
Appendix~\ref{app:refinements} & Calibrated geometric and noise refinements & PROVED\\
Theorems~\ref{thm:norm_diameter}--\ref{thm:all_norm_rates} & Full-data length for every norm budget & PROVED\\
Theorem~\ref{thm:all_norm_mse} & Uniform norm-budget squared risk & PROVED\\
Theorem~\ref{thm:weighted_bounds} & Orthogonal weighted upper/lower envelopes & PROVED\\
Proposition~\ref{prop:flat_loading} & Sharp flat-support and rotated-loading rates & PROVED\\
Original Q2 & Coordinate lower scale for all $K\ge3$ & REFUTED\\
Q3 & Actual-rank adaptation and sharp general loading geometry & OPEN\\
Q4 & Polynomial-time minimax confidence intervals & OPEN\\
\bottomrule
\end{tabularx}
\caption{Updated ledger. The earlier cubic and worst-case statistical gaps are closed for fixed caps; the norm-budget restriction is removed by the new results. The low-degree result is not an unrestricted computational-hardness theorem.}\label{tab:status}
\end{table}

The remaining theory concerns efficient minimax interval construction, actual-rank adaptation, arbitrary loading geometry, sharp cap dependence and matching bounds for general unequal loadings. The norm/noise phase diagram is now characterized for equal and flat orthogonal loadings. The noise extensions require their stated calibration and information. No empirical result for exact minimax projection or real-data/LLM validation is asserted. The priority of the new sharp-rate and scalar low-degree statements requires independent literature scrutiny.
